\documentclass[11pt,a4paper]{article}

\usepackage[a4paper,margin=2.6cm]{geometry}
\usepackage[T1]{fontenc}
\usepackage[utf8]{inputenc}
\usepackage[french]{babel}
\usepackage{newtxtext,newtxmath}
\usepackage{microtype}
\usepackage{amsmath}
\usepackage{enumitem}
\usepackage{xcolor}
\usepackage{titlesec}
\usepackage{hyperref}

\definecolor{primary}{RGB}{30,50,82}
\definecolor{secondary}{RGB}{117,96,61}

\hypersetup{
    colorlinks=true,
    linkcolor=primary,
    pdftitle={Propositions de plans sur les attaques par injection de nœuds},
    pdfsubject={Plans de présentation sur les attaques et défenses des GNN}
}

\setlength{\parindent}{0pt}
\setlength{\parskip}{5pt}
\setlength{\emergencystretch}{2em}
\setlist[itemize]{leftmargin=1.5em,itemsep=3pt,topsep=3pt}

\titleformat{\section}
  {\Large\bfseries\color{primary}}
  {\thesection.}{0.5em}{}
  [\vspace{2pt}{\color{secondary}\titlerule[0.6pt]}]
\titleformat{\subsection}
  {\normalsize\sffamily\bfseries\color{primary}}
  {\thesubsection}{0.5em}{}

\newcommand{\duree}[1]{\hfill\textit{#1}}

\begin{document}

\begin{center}
    {\small\sffamily\scshape\color{secondary}Document de travail}\\[8pt]
    {\LARGE\bfseries\color{primary}Propositions de plans}\\[4pt]
    {\large Attaques par injection de nœuds et robustesse des GNN}\\[10pt]
    {\color{secondary}\rule{0.18\linewidth}{1pt}}\\[8pt]
    {\small Présentation orale : 15 minutes\quad\textbullet\quad Démonstration : 5 minutes}
\end{center}

\tableofcontents
\vspace{1em}

\section{Paradigmes d'optimisation}

\textbf{Fil conducteur.} Classer les attaques selon la façon dont elles résolvent
le problème d'optimisation, puis examiner les défenses et les adaptations des
attaquants.

\subsection{Introduction et modélisation des GIA}\duree{3 min}
\begin{itemize}
    \item Rappeler le passage de messages dans les réseaux de neurones sur graphes
    (GNN) et expliquer la vulnérabilité créée par l'agrégation des voisinages.
    \item Comparer l'injection de nœuds (GIA) à la modification du graphe (GMA),
    en soulignant le réalisme de l'injection et sa formalisation par blocs à
    l'aide de $A'$ et $X'$.
    \item Présenter les contraintes usuelles : budget de nœuds, degré maximal,
    et injection en une fois ou de manière séquentielle.
\end{itemize}

\subsection{Attaques fondées sur les gradients}\duree{3 min}
\begin{itemize}
    \item Exposer le problème de discrétisation : l'ajout d'une arête est une
    décision binaire, alors que les méthodes d'optimisation opèrent souvent
    dans un espace continu.
    \item Présenter les premières approches gloutonnes, dont Greedy et
    Greedy-GAN, ainsi que leur coût de calcul.
    \item Expliquer la linéarisation du modèle substitut et l'approximation
    rapide fondée sur le signe du gradient, avec FGSM adapté et AFGSM, pour
    traiter des graphes de grande taille.
\end{itemize}

\subsection{Topologie et apprentissage séquentiel}\duree{3 min}
\begin{itemize}
    \item Décrire la vulnérabilité topologique des nœuds de faible degré et
    l'exploitation de cette propriété pour sélectionner les arêtes d'injection
    dans TDGIA.
    \item Présenter l'optimisation sans dérivation directe par rapport au graphe,
    formulée comme un processus de décision markovien (MDP) et résolue par
    apprentissage par renforcement hiérarchique dans NIPA.
\end{itemize}

\subsection{Défenses contre les injections}\duree{3 min}
\begin{itemize}
    \item Introduire la rupture de l'homophilie locale comme faiblesse exploitable
    par les attaques par injection.
    \item Comparer les mécanismes de filtrage fondés sur l'homophilie ou
    l'attention (GNNGuard, EGNNGuard et RGAT), qui peuvent écarter les arêtes
    suspectes.
    \item Présenter les méthodes qui stabilisent les représentations internes :
    RobustGCN, LayerNorm et l'entraînement contradictoire FLAG.
\end{itemize}

\subsection{Synthèse critique et démonstration}\duree{3 min + 5 min de démonstration}
\begin{itemize}
    \item Montrer comment la régularisation d'homophilie HAO permet aux attaques
    de préserver leur discrétion et de contourner certaines défenses.
    \item Comparer les méthodes selon leur complexité et leur transférabilité
    entre architectures GNN.
    \item Conclure par une démonstration pratique d'injection et de défense.
\end{itemize}

\section{Anatomie d'une injection : connexions et attributs}

\textbf{Fil conducteur.} Décomposer une attaque en deux décisions techniques :
où connecter les nœuds injectés ($A_{\text{inj}}$) et quels attributs leur
attribuer ($X_{\text{inj}}$).

\subsection{Modèle de menace et propagation}\duree{3 min}
\begin{itemize}
    \item Expliquer la vulnérabilité structurelle des GNN et distinguer la
    modification du graphe (GMA) de l'injection de nœuds (GIA).
    \item Définir les cadres d'attaque : empoisonnement (\textit{poisoning}),
    notamment NIPA et AFGSM, ou évasion (\textit{evasion}), notamment TDGIA et
    HAO ; distinguer également les attaques ciblées des attaques globales.
\end{itemize}

\subsection{Choix des connexions : déterminer \texorpdfstring{$A_{\text{inj}}$}{A inj}}\duree{3 min}
\begin{itemize}
    \item Comparer la connexion directe à la cible et la connexion indirecte par
    l'intermédiaire de ses voisins.
    \item Présenter les stratégies fondées sur les gradients de perte : Greedy,
    AFGSM et MetaGIA.
    \item Examiner les stratégies heuristiques et séquentielles : score de
    vulnérabilité topologique dans TDGIA et sélection d'actions par renforcement
    dans NIPA.
\end{itemize}

\subsection{Choix des attributs : générer \texorpdfstring{$X_{\text{inj}}$}{X inj}}\duree{3 min}
\begin{itemize}
    \item Comparer l'optimisation d'attributs continus, par descente de gradient
    projetée (PGD) ou optimisation lisse dans TDGIA, aux choix discrets fondés
    sur le signe du gradient.
    \item Traiter le réalisme des attributs à l'aide de contraintes de
    cooccurrence (AFGSM), d'un discriminateur adversaire (Greedy-GAN) et de la
    préservation de l'homophilie (HAO).
\end{itemize}

\subsection{Détecter et neutraliser l'injection}\duree{3 min}
\begin{itemize}
    \item Expliquer pourquoi une injection naïve est plus facile à bloquer,
    notamment lorsqu'elle modifie la distribution d'homophilie.
    \item Décrire les défenses structurelles par élagage d'arêtes et seuil de
    similarité cosinus, comme GNNGuard et EGNNGuard.
    \item Comparer les défenses appliquées aux couches de convolution :
    normalisation (LayerNorm), distributions gaussiennes (RobustGCN) et
    architectures alternatives (SGCN, APPNP et TAGCN).
\end{itemize}

\subsection{Synthèse et démonstration}\duree{3 min + 5 min de démonstration}
\begin{itemize}
    \item Construire une matrice de synthèse comparant puissance d'attaque,
    coût algorithmique et furtivité face aux défenses.
    \item Présenter le protocole expérimental et lancer la démonstration.
\end{itemize}

\section{Évolution du domaine : la course aux armements}

\textbf{Fil conducteur.} Raconter l'évolution des méthodes comme une succession
d'attaques et de réponses défensives, depuis les premières attaques par
modification jusqu'aux attaques furtives.

\subsection{Genèse : pourquoi injecter plutôt que modifier ?}\duree{3 min}
\begin{itemize}
    \item Présenter les limites des premières attaques sur graphes, telles que
    Nettack, puis l'émergence des attaques par faux nœuds.
    \item Formaliser l'injection et expliquer l'effet théorique d'un nœud
    malveillant sur les couches de convolution d'un GNN.
\end{itemize}

\subsection{Première génération : discrétisation et passage à l'échelle}\duree{3 min}
\begin{itemize}
    \item Présenter les approches pionnières d'empoisonnement : algorithmes
    gloutons (Greedy-GAN), Q-learning hiérarchique (NIPA) et linéarisation
    rapide (AFGSM).
    \item Comparer leur complexité algorithmique, de $\mathcal{O}(n^2)$ à
    $\mathcal{O}(n)$ selon les méthodes et les hypothèses considérées.
\end{itemize}

\subsection{Deuxième génération : évasion à grande échelle}\duree{3 min}
\begin{itemize}
    \item Présenter le passage à des attaques en boîte noire au moment de
    l'inférence, sans réentraîner le GNN, illustré par le contexte du benchmark
    KDD Cup 2020.
    \item Expliquer TDGIA et son exploitation de la topologie pour maximiser
    l'impact avec une proportion réduite de nœuds injectés.
\end{itemize}

\subsection{Riposte : exploiter les failles des attaques}\duree{3 min}
\begin{itemize}
    \item Comparer l'effet des GIA et des GMA sur l'homophilie du graphe, à
    partir des résultats théoriques et empiriques.
    \item Présenter les défenses fondées sur le filtrage de l'homophilie
    (GNNGuard), la variance (RobustGCN), la normalisation (LayerNorm) et
    l'entraînement robuste.
\end{itemize}

\subsection{Troisième génération : furtivité, bilan et démonstration}\duree{3 min + 5 min de démonstration}
\begin{itemize}
    \item Expliquer comment l'objectif harmonieux HAO vise à préserver
    l'homophilie et à contourner les défenses.
    \item Présenter un tableau comparatif des méthodes, puis lancer la
    démonstration interactive.
\end{itemize}

\section{Propriétés des graphes : degré, homophilie, spectre et échelle}

\textbf{Fil conducteur.} Organiser l'exposé autour des propriétés structurelles
qui déterminent la vulnérabilité des nœuds, la discrétion des injections et
l'efficacité des défenses.

\subsection{Connectivité locale et modèle d'injection}\duree{3 min}
\begin{itemize}
    \item Montrer comment la topologie locale influe sur les prédictions d'un
    GNN.
    \item Formaliser l'injection de nœuds et les budgets imposés à la structure
    du graphe.
\end{itemize}

\subsection{Degré et centralité dans les attaques}\duree{3 min}
\begin{itemize}
    \item Analyser la vulnérabilité des nœuds en fonction de leur degré et
    expliquer pourquoi les nœuds de faible degré peuvent être plus sensibles.
    \item Examiner l'influence de la normalisation de la matrice d'adjacence,
    notamment la normalisation par ligne et la normalisation symétrique
    $\tilde{D}^{-1/2}\tilde{A}\tilde{D}^{-1/2}$.
    \item Relier ces propriétés aux méthodes TDGIA et AFGSM.
\end{itemize}

\subsection{Statistiques globales et optimisation séquentielle}\duree{3 min}
\begin{itemize}
    \item Étudier les moyens de limiter les changements à la loi de puissance
    des degrés, au nombre de composantes connexes et au nombre de triangles,
    notamment dans NIPA et Greedy-GAN.
    \item Comparer l'optimisation globale à l'ajout séquentiel des nœuds
    injectés.
\end{itemize}

\subsection{Homophilie et mécanismes de défense}\duree{3 min}
\begin{itemize}
    \item Définir l'homophilie centrée sur les nœuds et l'homophilie centrée
    sur les arêtes.
    \item Présenter les défenses fondées sur la restauration de l'homophilie
    (GNNGuard et élagage par similarité cosinus) ainsi que la régularisation
    des représentations (LayerNorm et RobustGCN).
    \item Examiner le contournement de ces défenses par la contrainte
    d'homophilie HAO.
\end{itemize}

\subsection{Synthèse multicritère et démonstration}\duree{3 min + 5 min de démonstration}
\begin{itemize}
    \item Récapituler les compromis entre attaques et défenses, sur des graphes
    de petite et de grande taille.
    \item Présenter et lancer une démonstration pratique sur le jeu de données
    Cora.
\end{itemize}

\end{document}
